\scanpage{59}
% l. 1
Cas $n=1$, sur corps \inserted{$k$} car. 0, $G_{0}$ fini

% ll. 2–3
\[
\begin{array}{cccl}
\mu_0,&\mu_1,&\beta&\text{invariant}\\[-2pt]
\shortparallel&\shortparallel&\shortparallel&\\
\lambda_0-1&\lambda_1-1&\alpha+2&
\end{array}
\]

% ll. 4–5
1) $\mu_0,\mu_1\ne0$ i.e. $\lambda_0,\lambda_1\ne1$ \quad ($\tau_0,\tau_1$ d'ordre fini)

$\lambda_0,\lambda_1$ racines de l'unité (d'ordre $\underline{\nu}_0,\underline{\nu}_1\unsure{\geqslant|>}2$)

% ll. 6–10
2) $\langle\pi,\pi'\rangle=\underset{\shortparallel\,\inserted{\deleted{\unsure{\beta}}\mu_0\mu_1-\beta}}{\rho}\ne0$ \quad i.e. $\unsure{E|\mathcal E}\simeq V\oplus k\pi$ \qquad On peut se borner : $V$

Système de deux ps. réflexions strictes $\tau_0,\tau_1$ dans $V$, avec conditions de non dégénérescence :

\[
\begin{aligned}
\text{a)}\;&\langle a_0,a'_1\rangle\ne0\\
\text{b)}\;&a_0,a_1\quad\text{lin. indép.}
\end{aligned}
\]

% ll. 11–15
I) La représentation dans $C=\check{\mathbb P}(V)$ admet un div. de degré 2 invariant

Il ne peut être réduit à un pt, car s'il y a une droite \inserted{$\subset V$} \uline{invariante} par $G_0$, il y a une \unsure{suppl|supplém.}, donc $G_0$ est \unsure{comm|commutatif} $\subset k^\times\times k^\times$, absurde.

% ll. 16–21
Donc $V=D\oplus D'$, et pour la $=$ raison que dessus, $D,D'$ ne sont pas invariantes par $\tau_0,\tau_1$ à la fois

\deleted{a) $\tau_0$ invarie $D,D'$, $\tau_1$ \unsure{non}} \inserted{Si $\tau_i$ invarie $D$,} c'est que ou bien $a'_i$ est nul sur $D$, ou bien $a^{\deleted{\unsure{\prime}}}_i\in D$. Dans le premier cas, on aurait $a_i\in D'$, dans le deuxième $a'_i$ nul sur $D'$. Ainsi, OPS $a_i\in D$, $a'_i$ nul sur $D'$, donc
\[
D=ka_i,\quad D'=\operatorname{Ker}a'_i
\]

% ll. 22–26
Considérons alors \deleted{$j\ne i\quad \unsure{0\leq j\leq1}$} \inserted{$j\ne i,\ j\in\{0,1\}$}, alors $\tau_jD\subset D'$, $\tau_jD'\subset D$,
\[
a_i+\underbrace{\langle a_i,a'_j\rangle}_{\deleted{\unsure{\mu_0\beta}}}a_j\in D'
\quad\text{i.e.}\quad
\underbrace{\langle a_i,a'_i\rangle}_{\mu_i}
+\underbrace{\langle a_i,a'_j\rangle\langle a_j,a'_i\rangle}_{\beta}=0
\]
d'où $\beta\ne0$, donc quitte à intervertir $\tau_0,\tau_1$ SPG que $i=0$, $j=1$ \quad $\beta=-\mu_0$. Exprimons aussi $\tau_1D'\subset D$ i.e. $\tau_1^2D\subset D$, $\tau_1^2=\mathrm{id}+(2+\mu_1)a'_1\otimes a_1$
% archive sheet number: 58
